% Abstract
Modern Site Reliability Engineering (SRE) teams have mature anomaly detectors; the bottleneck is no longer ``is something wrong?'' but ``what is wrong, has it happened before, and what should we do?''. We construct and characterize a memory-augmented incident-triage system that, given a 5-minute window of telemetry (logs, traces, metrics, Kubernetes events), produces three outputs simultaneously: (i) a calibrated triage probability, (ii) a ranked top-5 list of past Jira tickets, and (iii) a novelty flag indicating whether the failure mode appears genuinely new.

Because no public dataset jointly contains rich telemetry and linked Jira tickets, we built our own. We forked Google's Online Boutique microservices demo, added production-realistic OpenTelemetry instrumentation across all ten services (per-RPC structured logs, dependency-boundary error logs, manual child spans, RED metrics) under a strict ``no scenario-label leakage'' policy, defined twenty-seven incident scenario families, ran one hundred independent fault-injection runs producing 6{,}720 labeled telemetry windows, and generated a humanized 347-ticket Jira memory corpus plus a 110-ticket distractor pool. All artifacts are reproducible from scripts and a frozen split manifest.

Over the project's lifetime we evolved through four method generations. Single-pipeline baselines (gradient-boosted classifier on numeric features, BM25 retrieval, log-template miner) each saturated on one axis but no single pipeline dominated the metric space. Seven trained pipelines---\hgb, fine-tuned bi-encoder retrieval, hybrid RRF fusion with rule-extracted and LLM-extracted knowledge graphs, log-sequence encoder, structured graph retrieval, and a three-stage \agent{} powered by Qwen~3.6 35B-A3B---each won on one axis while losing on others. This Pareto-incomparable structure motivated our central contribution: the Tiered Cascade Hybrid (\tch{}), a four-layer composition that uses the right pipeline for the right sub-task. L1 stacks six triage scores via class-balanced logistic regression. L2 fuses four retrievers via Reciprocal Rank Fusion with an overlap-rerank step for top-1 precision. L3 detects novel failure modes via three OR-combined signals (LLM agent verification, retrieval-confidence threshold, and a learned per-window classifier added in late-stage refinement). L4 composes the final output. The cascade runs at 16 microseconds per window after caching---essentially free at inference.

On 1{,}008 in-distribution test windows, \tch{} achieves Hit@5 = 0.912, Hit@1 = 0.722, MRR = 0.794, strict PR-AUC = 0.9998, and novel-recall = 0.793 at novel-precision = 0.940. The headline numbers match or beat every single-pipeline baseline across all axes simultaneously; the novel-recall figure represents a +388\% relative improvement over the locked v2f baseline's 0.163 at preserved precision, and is the largest single improvement of the project. We document eight late-stage refinement experiments (G1--G8): three improved the cascade and three failed in informative ways, with the failures forming a consistent pattern---any single-component intervention that tries to outscore the cascade's existing multi-retriever consensus underperforms it. The cascade is robust to memory noise (Hit@5 drops only 2\% relative at 50\% distractor ratio) and to family-distribution shift (novelty F1 within 11\% relative under leave-one-family-out cross-validation).

This report covers the full timeline: how we collected the data, what we instrumented in the microservices, which methods we tried and why, how the hybrid emerged from the constraint that no single model dominated, the cascade's mathematical specification, and what eight late-stage experiments taught us. It is intended as a comprehensive technical record suitable for institutional review and as the basis of a future peer-reviewed submission.
